% Standalone research note for the hybrid-local-solver project.
%
% The note proves that degree-normalized PPR is maximized at its seed, fixes
% the resulting response-row reciprocity factor, and derives the general-graph
% terminal potential bound for nonnegative residual/output maps.

\documentclass[11pt]{article}

% Common class-neutral preamble for standalone notes under manuscript/notes/.
% Note entry points all live exactly two directories below manuscript/.

\usepackage[margin=1in]{geometry}
\usepackage{amsmath,amssymb,amsthm,mathtools,bm}
\usepackage{algorithm}
\usepackage{algorithmic}
\usepackage{booktabs,tabularx,multirow,array,longtable}
\usepackage{enumitem}
\usepackage{microtype}
\usepackage[numbers,sort&compress]{natbib}
\usepackage{xcolor}
\usepackage[colorlinks=true,citecolor=blue,linkcolor=black,urlcolor=blue]{hyperref}
\usepackage[capitalize,nameinlink,noabbrev]{cleveref}

\newtheorem{theorem}{Theorem}[section]
\newtheorem{lemma}[theorem]{Lemma}
\newtheorem{proposition}[theorem]{Proposition}
\newtheorem{corollary}[theorem]{Corollary}
\newtheorem{conjecture}[theorem]{Conjecture}
\theoremstyle{definition}
\newtheorem{definition}[theorem]{Definition}
\newtheorem{assumption}[theorem]{Assumption}
\newtheorem{problem}[theorem]{Problem}
\newtheorem{observation}[theorem]{Observation}
\theoremstyle{remark}
\newtheorem{remark}[theorem]{Remark}
\theoremstyle{plain}

% Canonical mathematical notation for the active manuscript.
%
% This file preserves the recurring notation style used in the author's
% NeurIPS 2024 and 2025 papers while excluding unrelated tensor and
% machine-learning boilerplate from those archived command collections.
%
% Prefix convention:
%   \v...  bold vectors
%   \m...  bold matrices
%   \g...  calligraphic sets, graphs, and families
%   \s...  blackboard-bold spaces and number systems

\newcommand{\method}{\textsc{HybridLocalSolver}}

% Font and scalar shorthands retained from the archived manuscripts.
\newcommand{\mc}{\mathcal}
\newcommand{\R}{\mathbb{R}}
\newcommand{\E}{\mathbb{E}}
\newcommand{\G}{\mathbb{G}}
\newcommand{\N}{\mathcal{N}}
% Accuracy namespaces are semantic: do not add a bare \eps alias.
% Degree-normalized residual tolerance of classical APPR is kept distinct from
% the objective-gap and proximal fixed-point tolerances of the RPPR sections.
\newcommand{\epsappr}{\varepsilon_{\mathrm{appr}}}
\newcommand{\epsobj}{\varepsilon_{\mathrm{obj}}}
\newcommand{\epspg}{\varepsilon_{\mathrm{pg}}}
\newcommand{\epsppr}{\varepsilon_{\mathrm{ppr}}}
\newcommand{\epsin}{\varepsilon_{\mathrm{in}}}
\newcommand{\epskkt}{\varepsilon_{\mathrm{kkt}}}
\newcommand{\epsburn}{\varepsilon_{\mathrm{burn}}}
\newcommand{\epssol}{\varepsilon_{\mathrm{sol}}}
\newcommand{\epspprinst}[1]{\varepsilon_{\mathrm{ppr},#1}}
\newcommand{\trans}{\mathsf{T}}
\newcommand{\grad}{\nabla}

% Delimiters and generic helpers.
% Norms and inner products use fixed-size delimiters by default.
% Use an optional size (e.g. \norm[\big]{...}) or * for automatic sizing.
\DeclarePairedDelimiter{\norm}{\lVert}{\rVert}
\newcommand{\abs}[1]{\left\lvert #1 \right\rvert}
\DeclarePairedDelimiterX{\inner}[2]{\langle}{\rangle}{#1, #2}
\newcommand{\ip}[2]{\inner{#1}{#2}}
\newcommand{\card}[1]{\left\lvert #1 \right\rvert}
\newcommand{\set}[1]{\left\{#1\right\}}
\newcommand{\ceil}[1]{\left\lceil #1 \right\rceil}
\newcommand{\floor}[1]{\left\lfloor #1 \right\rfloor}
\newcommand{\comp}[1]{\overline{#1}}
\newcommand{\conj}[1]{\overline{#1}}
\newcommand{\mean}[1]{\overline{#1}}
\newcommand{\bigO}{\mathcal{O}}
\newcommand{\softO}{\widetilde{\mathcal{O}}}
\newcommand{\softOmega}{\widetilde{\Omega}}
\newcommand{\pos}[1]{\left[#1\right]_{+}}

% Bold lowercase vectors.
\newcommand{\va}{\bm{a}}
\newcommand{\vb}{\bm{b}}
\newcommand{\vc}{\bm{c}}
\newcommand{\vd}{\bm{d}}
\newcommand{\ve}{\bm{e}}
\newcommand{\vf}{\bm{f}}
\newcommand{\vg}{\bm{g}}
\newcommand{\vh}{\bm{h}}
\newcommand{\vi}{\bm{i}}
\newcommand{\vj}{\bm{j}}
\newcommand{\vk}{\bm{k}}
\newcommand{\vl}{\bm{l}}
\newcommand{\vm}{\bm{m}}
\newcommand{\vn}{\bm{n}}
\newcommand{\vo}{\bm{o}}
\newcommand{\vp}{\bm{p}}
\newcommand{\vq}{\bm{q}}
\newcommand{\vr}{\bm{r}}
\newcommand{\vs}{\bm{s}}
\newcommand{\vt}{\bm{t}}
\newcommand{\vu}{\bm{u}}
\newcommand{\vv}{\bm{v}}
\newcommand{\vw}{\bm{w}}
\newcommand{\vx}{\bm{x}}
\newcommand{\vy}{\bm{y}}
\newcommand{\vz}{\bm{z}}

% Bold Greek vectors and special vectors.
\newcommand{\valpha}{\bm{\alpha}}
\newcommand{\vbeta}{\bm{\beta}}
\newcommand{\vdelta}{\bm{\delta}}
\newcommand{\vepsilon}{\bm{\epsilon}}
\newcommand{\veta}{\bm{\eta}}
\newcommand{\vgamma}{\bm{\gamma}}
\newcommand{\vlambda}{\bm{\lambda}}
\newcommand{\vmu}{\bm{\mu}}
\newcommand{\vomega}{\bm{\omega}}
\newcommand{\vphi}{\bm{\phi}}
\newcommand{\vpi}{\bm{\pi}}
\newcommand{\vpsi}{\bm{\psi}}
\newcommand{\vrho}{\bm{\rho}}
\newcommand{\vsigma}{\bm{\sigma}}
\newcommand{\vtheta}{\bm{\theta}}
\newcommand{\vxi}{\bm{\xi}}
\newcommand{\vell}{\bm{\ell}}
\newcommand{\vDelta}{\bm{\Delta}}
\newcommand{\vGamma}{\bm{\Gamma}}
\newcommand{\vLambda}{\bm{\Lambda}}
\newcommand{\vPhi}{\bm{\Phi}}
\newcommand{\vSigma}{\bm{\Sigma}}
\newcommand{\vzero}{\bm{0}}
\newcommand{\vone}{\bm{1}}
\newcommand{\one}{\mathbf{1}}
\newcommand{\zero}{\vzero}
\newcommand{\eunit}[1]{\bm{e}_{#1}}
\newcommand{\vDeltax}{\bm{\Delta}\bm{x}}

% Bold uppercase matrices.
\newcommand{\mA}{\bm{A}}
\newcommand{\mB}{\bm{B}}
\newcommand{\mC}{\bm{C}}
\newcommand{\mD}{\bm{D}}
\newcommand{\mE}{\bm{E}}
\newcommand{\mF}{\bm{F}}
\newcommand{\mG}{\bm{G}}
\newcommand{\mH}{\bm{H}}
\newcommand{\mI}{\bm{I}}
\newcommand{\mJ}{\bm{J}}
\newcommand{\mK}{\bm{K}}
\newcommand{\mL}{\bm{L}}
\newcommand{\mM}{\bm{M}}
\newcommand{\mN}{\bm{N}}
\newcommand{\mO}{\bm{O}}
\newcommand{\mP}{\bm{P}}
\newcommand{\mQ}{\bm{Q}}
\newcommand{\mR}{\bm{R}}
\newcommand{\mS}{\bm{S}}
\newcommand{\mT}{\bm{T}}
\newcommand{\mU}{\bm{U}}
\newcommand{\mV}{\bm{V}}
\newcommand{\mW}{\bm{W}}
\newcommand{\mX}{\bm{X}}
\newcommand{\mY}{\bm{Y}}
\newcommand{\mZ}{\bm{Z}}

% Bold Greek matrices retained from the graph papers.
\newcommand{\mBeta}{\bm{\beta}}
\newcommand{\mDelta}{\bm{\Delta}}
\newcommand{\mGamma}{\bm{\Gamma}}
\newcommand{\mLambda}{\bm{\Lambda}}
\newcommand{\mOmega}{\bm{\Omega}}
\newcommand{\mPhi}{\bm{\Phi}}
\newcommand{\mPi}{\bm{\Pi}}
\newcommand{\mSigma}{\bm{\Sigma}}
\newcommand{\mTheta}{\bm{\Theta}}

% Calligraphic symbols.
\newcommand{\gA}{\mathcal{A}}
\newcommand{\gB}{\mathcal{B}}
\newcommand{\gC}{\mathcal{C}}
\newcommand{\gD}{\mathcal{D}}
\newcommand{\gE}{\mathcal{E}}
\newcommand{\gF}{\mathcal{F}}
\newcommand{\gG}{\mathcal{G}}
\newcommand{\gH}{\mathcal{H}}
\newcommand{\gI}{\mathcal{I}}
\newcommand{\gJ}{\mathcal{J}}
\newcommand{\gK}{\mathcal{K}}
\newcommand{\gL}{\mathcal{L}}
\newcommand{\gM}{\mathcal{M}}
\newcommand{\gN}{\mathcal{N}}
\newcommand{\gO}{\mathcal{O}}
\newcommand{\gP}{\mathcal{P}}
\newcommand{\gQ}{\mathcal{Q}}
\newcommand{\gR}{\mathcal{R}}
\newcommand{\gS}{\mathcal{S}}
\newcommand{\gT}{\mathcal{T}}
\newcommand{\gU}{\mathcal{U}}
\newcommand{\gV}{\mathcal{V}}
\newcommand{\gW}{\mathcal{W}}
\newcommand{\gX}{\mathcal{X}}
\newcommand{\gY}{\mathcal{Y}}
\newcommand{\gZ}{\mathcal{Z}}

% Blackboard-bold symbols.
\newcommand{\sA}{\mathbb{A}}
\newcommand{\sB}{\mathbb{B}}
\newcommand{\sC}{\mathbb{C}}
\newcommand{\sD}{\mathbb{D}}
\newcommand{\sE}{\mathbb{E}}
\newcommand{\sF}{\mathbb{F}}
\newcommand{\sG}{\mathbb{G}}
\newcommand{\sH}{\mathbb{H}}
\newcommand{\sI}{\mathbb{I}}
\newcommand{\sJ}{\mathbb{J}}
\newcommand{\sK}{\mathbb{K}}
\newcommand{\sL}{\mathbb{L}}
\newcommand{\sM}{\mathbb{M}}
\newcommand{\sN}{\mathbb{N}}
\newcommand{\sO}{\mathbb{O}}
\newcommand{\sP}{\mathbb{P}}
\newcommand{\sQ}{\mathbb{Q}}
\newcommand{\sR}{\mathbb{R}}
\newcommand{\sS}{\mathbb{S}}
\newcommand{\sT}{\mathbb{T}}
\newcommand{\sU}{\mathbb{U}}
\newcommand{\sV}{\mathbb{V}}
\newcommand{\sW}{\mathbb{W}}
\newcommand{\sX}{\mathbb{X}}
\newcommand{\sY}{\mathbb{Y}}
\newcommand{\sZ}{\mathbb{Z}}

% Frequently used operators.
\DeclareMathOperator{\Cov}{Cov}
\DeclareMathOperator{\Var}{Var}
\DeclareMathOperator{\diag}{diag}
\DeclareMathOperator{\nei}{Nei}
\DeclareMathOperator{\nnz}{nnz}
\DeclareMathOperator{\pr}{pr}
\DeclareMathOperator{\prox}{prox}
\DeclareMathOperator{\push}{push}
\DeclareMathOperator{\res}{res}
\DeclareMathOperator{\rel}{Rel}
\DeclareMathOperator{\relu}{ReLU}
\DeclareMathOperator{\sign}{sign}
\DeclareMathOperator{\supp}{supp}
\DeclareMathOperator{\tr}{tr}
\DeclareMathOperator{\Tr}{Tr}
\DeclareMathOperator{\vol}{vol}
\DeclareMathOperator{\work}{work}
\DeclareMathOperator{\Work}{Work}
\DeclareMathOperator*{\argmax}{arg\,max}
\DeclareMathOperator*{\argmin}{arg\,min}

% Project-wide names and claim-status labels used by standalone research notes.
% Mathematical notation belongs in math_commands.tex.

% Algorithms and solver families.
\newcommand{\AESP}{\textsc{AESP}}
\newcommand{\APPR}{\textsc{APPR}}
\newcommand{\ASPR}{\textsc{ASPR}}
\newcommand{\FISTA}{\textsc{FISTA}}
\newcommand{\ISTA}{\textsc{ISTA}}
\newcommand{\LOCAPPR}{\textsc{LocAPPR}}
\newcommand{\LOCGD}{\textsc{LocGD}}
\newcommand{\LOCSOR}{\textsc{LocSOR}}
\newcommand{\LOCCH}{\textsc{LocCH}}
\newcommand{\LOCHB}{\textsc{LocHB}}
\newcommand{\LocProxGS}{\textsc{LocProxGS}}
\newcommand{\Hybrid}{\textsc{AESP--LocSOR}}

% Claim classifications required by the repository research-note protocol.
\newcommand{\statussource}{\textbf{Source result}}
\newcommand{\statusproved}{\textbf{Proved here}}
\newcommand{\statusconditional}{\textbf{Conditional}}
\newcommand{\statusconjecture}{\textbf{Open / conjectural}}
\newcommand{\statusopen}{\textbf{Open}}
\newcommand{\statusempirical}{\textbf{Empirical}}
\newcommand{\statusobservation}{\textbf{Empirical observation}}
\newcommand{\statusrefuted}{\textbf{Refuted route}}

% Compatibility names used by the older complete-note presentation layer.
\newcommand{\Status}[2]{\textbf{#1:} #2}
\newcommand{\SourceStatus}{\textnormal{\textsc{Source result}}}
\newcommand{\ProvedStatus}{\textnormal{\textsc{Proved here}}}
\newcommand{\ConditionalStatus}{\textnormal{\textsc{Conditional}}}
\newcommand{\ComputedStatus}{\textnormal{\textsc{Computed}}}
\newcommand{\RefutedStatus}{\textnormal{\textsc{Refuted route}}}
\newcommand{\OpenStatus}{\textnormal{\textsc{Open}}}



\title{A Seed Maximum Principle for Personalized PageRank\\
and a General Output-Map Potential Bound}
\author{Baojian Zhou\\
\small Working research note for \texttt{hybrid-local-solver}}
\date{August 2026}

\begin{document}
\maketitle

\begin{abstract}
For the source-aligned personalized PageRank vector $\pi$ seeded at $v$, we
prove the discrete maximum principle
\[
    \frac{\pi_u}{d_u}\le \frac{\pi_v}{d_v}
    \qquad (u\in V).
\]
Combining it with reversibility gives the exact response-row formula
\[
 \frac{\operatorname{pr}(e_u)_v}{\pi_v}
 =\frac{d_v\pi_u}{d_u\pi_v}
 =\frac{\pi_u/d_u}{\pi_v/d_v}\le1.
\]
Thus the response-row constant relevant to the monotone unsettled-value
potential equals one on every connected graph, not merely on the centre-seeded
star.  This corrects an inverted degree ratio in an earlier proof-campaign
transcription.

As a consequence, if a one-hop state has nonnegative residual $r$ and reports
$z+Mr$ using any nonnegative output map with column mass at most $B$, then a
degree-normalized $\epsppr$ solution implies
\[
 \frac{(H^{-1}r)_v}{\pi_v}
 \le \frac{\|r\|_1}{\gamma_{\alpha}}
 \le \frac{\epsppr\vol(V)}{1-\gamma_{\alpha} B}
 \qquad(\gamma_{\alpha} B<1).
\]
This closes the general-graph terminal-potential question and extends the
associated monotone seed-operation lower bound to output maps.  It does not
apply to signed residuals, signed output maps, or stronger response
primitives, and it does not resolve the signed-relaxation lower-bound problem.
\end{abstract}

% Canonical source-aligned PageRank/RPPR reference formulation.
%
% This fragment is intentionally heading-free at the section level so that the
% active paper and every standalone note can input it. It is a manuscript
% reference convention, not an implementation-wide residual decision.
% Primary source: Fountoulakis and Martinez-Rubio, arXiv:2602.21138v2,
% PDF p. 3, Section 3 and equation (RPPR).

\subsection{Common source-aligned PageRank and RPPR model}
\label{subsec:shared-source-aligned-problem}

All documents in this manuscript workspace use the following reference
language. It follows the plain italic notation of the comparison paper:
\(x,s,A,D,Q,I\) denote column vectors or matrices as determined by their
displayed domains. This is the scoped exception to the project's usual
bold-vector and bold-matrix typography. A note may impose a stronger
assumption after this block, but it must not redefine these objects.

Let \(G=(V,E)\) be a finite, simple, undirected, unweighted graph without
isolated vertices, let \(V=[n]:=\{1,\ldots,n\}\), and let
\(A\in\{0,1\}^{n\times n}\) be its symmetric adjacency matrix. For
\(i\in V\), write
\[
    \mathcal N(i):=\{j\in V:\{i,j\}\in E\},
    \qquad
    d_i:=|\mathcal N(i)|,
    \qquad
    D:=\diag(d_1,\ldots,d_n).
\]
For \(S\subseteq V\), define
\[
    \mathcal N(S):=\bigcup_{i\in S}\mathcal N(i),
    \qquad
    \partial S:=\mathcal N(S)\setminus S,
    \qquad
    \operatorname{Ext}(S):=V\setminus(S\cup\partial S),
\]
\begin{equation}
    \vol(S):=\sum_{i\in S}d_i.
    \label{eq:shared-volume}
\end{equation}
The support of a vector is
\(\supp(x):=\{i\in[n]:x_i\neq0\}\). Matrix inequalities use the Loewner
order, and all unqualified vector inequalities are coordinatewise.

The seed is a column distribution
\begin{equation}
    s\in\R_+^n,
    \qquad
    \inner{\one}{s}=1.
    \label{eq:shared-seed}
\end{equation}
The single-seed specialization is always written \(s=e_v\), where \(v\in V\)
is the seed vertex. Thus \(s\) is never used as a vertex label. For
\(\alpha\in(0,1]\), define
\begin{equation}
\begin{aligned}
    \mathcal L&:=I-D^{-1/2}AD^{-1/2},\\
    Q&:=\alpha I+\frac{1-\alpha}{2}\mathcal L
      =\frac{1+\alpha}{2}I
       -\frac{1-\alpha}{2}D^{-1/2}AD^{-1/2},\\
    b&:=\alpha D^{-1/2}s.
\end{aligned}
    \label{eq:shared-pagerank-matrices}
\end{equation}
Since \(0\preceq\mathcal L\preceq2I\),
\(\alpha I\preceq Q\preceq I\). The shared smooth PageRank quadratic is
\begin{equation}
    f(x):=\frac12\inner{x}{Qx}-\inner{b}{x},
    \qquad
    \nabla f(x)=Qx-b.
    \label{eq:shared-pagerank-objective}
\end{equation}
Its unique unregularized minimizer and associated lazy personalized PageRank
vector are
\begin{equation}
    x^0:=Q^{-1}b,
    \qquad
    \pi:=D^{1/2}x^0.
    \label{eq:shared-ppr-solution}
\end{equation}

For a regularization parameter \(\rho>0\), reserve \(g_\rho\) for the
nonsmooth weighted \(\ell_1\) term and define
\begin{equation}
    g_\rho(x):=\alpha\rho\norm{D^{1/2}x}_1,
    \qquad
    F_\rho(x):=f(x)+g_\rho(x),
    \qquad
    x^\star(\rho):=\argmin_{x\in\R^n}F_\rho(x).
    \label{eq:shared-rppr-objective}
\end{equation}
When \(\rho\) is fixed, \(x^\star\) abbreviates \(x^\star(\rho)\), never the
unregularized point \(x^0\). Write
\[
    S^\star(\rho):=\supp(x^\star(\rho)),
    \qquad
    I^\star(\rho):=[n]\setminus S^\star(\rho).
\]
The source records \(x^\star(\rho)\geq0\),
\(\vol(S^\star(\rho))\leq1/\rho\), and the coordinatewise conditions
\begin{equation}
    \nabla_i f(x^\star(\rho))
    \in
    \begin{cases}
        \{-\alpha\rho\sqrt{d_i}\}, & x_i^\star(\rho)>0,\\
        [-\alpha\rho\sqrt{d_i},0], & x_i^\star(\rho)=0.
    \end{cases}
    \label{eq:shared-rppr-kkt}
\end{equation}

The common computational unit is one repeated adjacency-list scan: scanning
coordinate \(i\) costs \(d_i\), and scanning a set \(S\) costs \(\vol(S)\).
Iteration count and wall-clock time do not replace this degree-weighted work
measure.

Finally, accuracy symbols remain distinct. The APPR threshold is
\(\epsappr\); \(\epsobj\) denotes an objective-gap target;
\(\epspg\) denotes the source experiment's proximal fixed-point tolerance; and
\(\epsppr\) may denote a document-scoped degree-normalized PPR error target.
No equality or conversion among these quantities is assumed. In particular,
this shared model does not resolve the repository-wide residual decision.


\section{Scope and push-scale formulation}
\label{sec:seed-scope}

\paragraph{Claim status.}
All mathematical statements in this note are \statusproved{} in the shared
exact-arithmetic model.  The note imports no theorem beyond the shared PPR
definition.  Cross-direction integration is \statusopen{} because the note
whose transcription motivated the correction was separately owned and was
left untouched.

Specialize the shared model to a single seed $s=e_v$ and put
\begin{equation}
 c_{\alpha}:=\frac{1-\alpha}{1+\alpha},
 \qquad
 \gamma_{\alpha}:=\frac{2\alpha}{1+\alpha}=1-c_{\alpha},
 \qquad
 H:=I-c_{\alpha} AD^{-1}.
 \label{eq:seed-push-operator}
\end{equation}
Then $H\pi=\gamma_{\alpha} e_v$.  For any $h\in\mathbb R^V$, define the discounted
response
\begin{equation}
 \operatorname{pr}(h):=\gamma_{\alpha} H^{-1}h,
 \qquad \operatorname{pr}(e_v)=\pi.
 \label{eq:seed-response}
\end{equation}
Since $c_{\alpha} AD^{-1}$ is nonnegative with column sums $c_{\alpha}<1$, its
Neumann series shows that $H^{-1}>0$ entrywise on a connected graph.

The state used in the output-map corollary is
\begin{equation}
 r:=\gamma_{\alpha} e_v-Hz,
 \qquad e:=\pi-z=H^{-1}r.
 \label{eq:seed-state}
\end{equation}
There $r\ge0$ is an explicit assumption.  The reported vector is $z+Mr$,
where $M\ge0$ and every column sum of $M$ is at most $B$.  The proof does not
need the additional one-hop sparsity of $M$, so it applies in particular to
the project output-map class.

\section{Reciprocity and the seed maximum principle}
\label{sec:seed-maximum}

\begin{lemma}[Discounted Green-kernel reciprocity; \statusproved]
\label{lem:seed-reciprocity}
For every $u,w\in V$,
\begin{equation}
 d_u(H^{-1})_{wu}=d_w(H^{-1})_{uw}.
 \label{eq:seed-reciprocity}
\end{equation}
Consequently, for PPR seeded at $v$,
\begin{equation}
 \frac{\operatorname{pr}(e_u)_v}{\pi_v}
 =\frac{d_v\pi_u}{d_u\pi_v}
 =\frac{\pi_u/d_u}{\pi_v/d_v}.
 \label{eq:seed-response-ratio}
\end{equation}
\end{lemma}

\begin{proof}
The matrix
\[
 K:=D^{-1}H=D^{-1}-c_{\alpha} D^{-1}AD^{-1}
\]
is symmetric positive definite.  Since $H=DK$, one has
$H^{-1}=K^{-1}D^{-1}$.  Symmetry of $K^{-1}$ gives
$(H^{-1})_{wu}=(K^{-1})_{wu}/d_u
=(K^{-1})_{uw}/d_u=(d_w/d_u)(H^{-1})_{uw}$, proving
\eqref{eq:seed-reciprocity}.  Multiply by $\gamma_{\alpha}$ and use
$\pi_u=\gamma_{\alpha}(H^{-1})_{uv}$ and
$\pi_v=\gamma_{\alpha}(H^{-1})_{vv}$ to obtain
\eqref{eq:seed-response-ratio}.
\end{proof}

\begin{theorem}[Degree-normalized seed maximum; \statusproved]
\label{thm:seed-normalized-maximum}
For every finite simple connected undirected unit-weight graph with at least
two vertices, every $\alpha\in(0,1)$, and every seed $v$,
\begin{equation}
 \max_{u\in V}\frac{\pi_u}{d_u}=\frac{\pi_v}{d_v},
 \label{eq:seed-normalized-maximum}
\end{equation}
and the maximum is strict away from $v$.
\end{theorem}

\begin{proof}
Put $q:=D^{-1}\pi$.  Multiplying $H\pi=\gamma_{\alpha} e_v$ by the identity
$\pi=Dq$ gives
\begin{equation}
 (D-c_{\alpha} A)q=\gamma_{\alpha} e_v.
 \label{eq:seed-harmonic-equation}
\end{equation}
The Neumann-series observation in \cref{sec:seed-scope} gives $q>0$.  At
every nonseed vertex $u\ne v$, \eqref{eq:seed-harmonic-equation} says
\begin{equation}
 q_u=c_{\alpha}\frac{1}{d_u}\sum_{w\sim u}q_w.
 \label{eq:seed-strict-subharmonic}
\end{equation}
If a nonseed vertex attained the positive global maximum $q_{\max}$, the
right side would be at most $c_{\alpha} q_{\max}<q_{\max}$, a contradiction.
Thus the maximum is attained only at $v$, which is exactly
\eqref{eq:seed-normalized-maximum}.
\end{proof}

\begin{corollary}[Exact response-row constant; \statusproved]
\label{cor:seed-response-row}
Define
\begin{equation}
 \Phi_v:=\max_{u\in V}\frac{\operatorname{pr}(e_u)_v}{\pi_v}.
 \label{eq:seed-phi-definition}
\end{equation}
Then
\begin{equation}
 \Phi_v
 =\max_{u\in V}\frac{d_v\pi_u}{d_u\pi_v}
 =1.
 \label{eq:seed-phi-one}
\end{equation}
\end{corollary}

\begin{proof}
Combine \cref{lem:seed-reciprocity,thm:seed-normalized-maximum}; equality
holds at $u=v$.
\end{proof}

\paragraph{Correction to the campaign transcription.}
The earlier formula
\[
 \frac{\operatorname{pr}(e_u)_v}{\pi_v}
 \stackrel{\mathrm{incorrect}}{=}
 \frac{d_u\pi_u}{d_v\pi_v}
\]
has the degree ratio reversed.  On an irregular graph the two expressions
differ.  The corrected ratio in \eqref{eq:seed-response-ratio} both matches
the star identity and turns the measured conjecture $\Phi_v=1$ into the
maximum principle above.

\section{Terminal potential and output maps}
\label{sec:seed-output-map}

\begin{lemma}[Nonnegative residual response; \statusproved]
\label{lem:seed-residual-response}
For every $r\ge0$,
\begin{equation}
 \Psi_v(r):=\frac{(H^{-1}r)_v}{\pi_v}
 \le \frac{\|r\|_1}{\gamma_{\alpha}}.
 \label{eq:seed-residual-response}
\end{equation}
\end{lemma}

\begin{proof}
Write $r=\gamma_{\alpha} r_{\rm mass}$.  By linearity and
\cref{cor:seed-response-row},
\[
 \Psi_v(r)
 =\sum_u\frac{\operatorname{pr}(e_u)_v}{\pi_v}(r_{\rm mass})_u
 \le\sum_u(r_{\rm mass})_u
 =\frac{\|r\|_1}{\gamma_{\alpha}}.
\]
\end{proof}

\begin{corollary}[General output-map terminal bound; \statusproved]
\label{cor:seed-output-terminal}
Suppose the state \eqref{eq:seed-state} has $r\ge0$.  Let $M\ge0$ have
column sums at most $B$, and suppose the reported output satisfies
\begin{equation}
 \max_u\frac{|\pi_u-z_u-(Mr)_u|}{d_u}\le\epsppr.
 \label{eq:seed-semantic-output}
\end{equation}
If $\gamma_{\alpha} B<1$, then
\begin{equation}
 \frac{\|r\|_1}{\gamma_{\alpha}}
 \le\frac{\epsppr\vol(V)}{1-\gamma_{\alpha} B},
 \qquad
 \Psi_v(r)
 \le\frac{\epsppr\vol(V)}{1-\gamma_{\alpha} B}.
 \label{eq:seed-output-terminal}
\end{equation}
\end{corollary}

\begin{proof}
Because $r\ge0$, the column-sum assumption gives
$\one^TMr\le B\|r\|_1$.  Also
$H^T\one=\gamma_{\alpha}\one$, and hence
\begin{equation}
 \one^TH^{-1}r=\frac{\|r\|_1}{\gamma_{\alpha}}.
 \label{eq:seed-error-mass}
\end{equation}
Summing the one-sided inequalities
$\pi_u-z_u-(Mr)_u\le\epsppr d_u$ from
\eqref{eq:seed-semantic-output} yields
\[
 \|r\|_1\left(\frac1{\gamma_{\alpha}}-B\right)
 \le\epsppr\vol(V).
\]
This is the first bound in \eqref{eq:seed-output-terminal}; the second follows
from \cref{lem:seed-residual-response}.
\end{proof}

\begin{corollary}[Monotone seed-operation count with output maps;
\statusproved]
\label{cor:seed-monotone-count}
Consider invariant-exact one-hop updates
\[
 z_u\mathrel{+}=\delta,
 \qquad r\mathrel{-}=\delta He_u,
\]
that preserve $r\ge0$, start from $z=0$, and charge $d_u$ per update.  Let
$N_v^+$ be the number of seed updates with $\delta>0$, put
$\kappa_v:=\gamma_{\alpha}/\pi_v$, and assume $\alpha\in(0,1)$.  Under the terminal
conditions of \cref{cor:seed-output-terminal}, if
$\epsppr\vol(V)<1-\gamma_{\alpha} B$, then
\begin{equation}
 N_v^+\ge
 \frac{\log\!\left((1-\gamma_{\alpha} B)/(\epsppr\vol(V))\right)}
      {-\log(1-\kappa_v)},
 \qquad
 \Work\ge d_vN_v^+.
 \label{eq:seed-monotone-count}
\end{equation}
\end{corollary}

\begin{proof}
The potential $\Psi=(\pi_v-z_v)/\pi_v$ starts at one and changes only at
seed updates, by $-\delta/\pi_v$.  Since $r\ge0$,
\[
 (H^{-1}r)_v\ge(H^{-1})_{vv}r_v
 =\frac{\pi_v}{\gamma_{\alpha}}r_v.
\]
Feasibility after a positive seed update gives $\delta\le r_v$, so that
$-\Delta\Psi\le\kappa_v\Psi$.  Negative seed updates increase $\Psi$, and
nonseed updates leave it unchanged.  Therefore
$\Psi_T\ge(1-\kappa_v)^{N_v^+}$.  The terminal upper bound in
\cref{cor:seed-output-terminal} is strictly below one under the displayed
assumption; combining the two bounds and taking logarithms proves
\eqref{eq:seed-monotone-count}.  Each counted update costs $d_v$.
\end{proof}

\section{Implications and limits}
\label{sec:seed-limits}

The response-row question is closed exactly: the correct constant is one on
every connected graph.  The proof also strengthens the terminal-potential
bound because the locality pattern of $M$ is unused; nonnegativity and column
mass are the only output-map properties required.

The result has four strict limits.
\begin{enumerate}
 \item It uses $r\ge0$ both in \cref{lem:seed-residual-response} and in the
   multiplicative step argument.  Signed-relaxation trajectories are outside
   the theorem.
 \item It uses $M\ge0$.  A signed response map can cancel mass and is not
   controlled by the column-sum argument.
 \item It is a class statement about invariant-exact one-hop updates.  It is
   not an information lower bound and says nothing against block elimination,
   Schur response, or another stronger primitive.
 \item It does not address whether a bounded-seed-degree graph family forces
   product-scale work for dissipative signed one-hop relaxation.  That is a
   separate spectral/work-distribution question.
\end{enumerate}

\end{document}
